\section*{Capítulo \ref{ch:g9:maps-of-the-earth} --- Mapas da Terra inteira}

\begin{solution}{exo:g9:maps-of-the-earth:1}
(a) Uma \omterm{def:g9:maps-of-the-earth:projection}{projeção cartográfica}. (b) Uma \omterm{def:g9:maps-of-the-earth:contour}{curva de nível}. (c) Uma \omterm{def:g9:maps-of-the-earth:numerical-scale}{escala
numérica}. (d) A \omterm{def:g9:maps-of-the-earth:contour}{equidistância}. (e) O \omterm{def:g9:maps-of-the-earth:remote-sensing}{sensoriamento remoto}.
\end{solution}

\begin{solution}{exo:g9:maps-of-the-earth:2}
(a) $4 \times 25\,000 = \num{100000}$~cm $= \qty{1}{km}$. (b) $3 \times
50\,000 = \num{150000}$~cm $= \qty{1.5}{km}$. (c) $2.5 \times 1\,000\,000
= \num{2500000}$~cm $= \qty{25}{km}$. (d) $\qty{1}{km} = \num{100000}$~cm;
$100\,000 \div 25\,000 = \qty{4}{cm}$.
\end{solution}

\begin{solution}{exo:g9:maps-of-the-earth:3}
(a) Falso: uma esfera não pode ser achatada sem deformar os comprimentos,
então todo \omterm{def:g3:continents-and-oceans:world-map}{mapa-múndi} deforma alguma coisa. (b) Verdadeiro. (c) Falso: um
\omterm{def:g2:drawing-to-map:map}{mapa} equivalente mantém as áreas, mas entorta as formas, sobretudo perto
das suas bordas. (d) Verdadeiro.
\end{solution}

\begin{solution}{exo:g9:maps-of-the-earth:4}
(a) A encosta é íngreme. (b) O terreno é suave. (c) Um morro, com o topo
dentro da menor linha. (d) Um ponto do terreno só tem uma \omterm{def:g4:world-physical-map:altitude}{altitude}, então
ele não pode ficar em duas linhas de \omterm{def:g4:world-physical-map:altitude}{altitudes} diferentes.
\end{solution}

\begin{solution}{exo:g9:maps-of-the-earth:5}
(a) A \ang{60}\,N, o \omterm{def:g1:landscapes-around-us:water}{lago} é desenhado 4 vezes grande demais: ele parece
tão grande quanto
$4 \times 1000 = \qty{4000}{km^2}$. No \omterm{def:g3:continents-and-oceans:equator}{equador}, as áreas não são
ampliadas, então o \omterm{def:g1:landscapes-around-us:water}{lago} desenhado do mesmo tamanho tem uma área
verdadeira de \qty{4000}{km^2}. (b) $4 \div 1.33 \approx 3$ vezes maior.
(c) As áreas, e as distâncias de norte a sul, crescem sem limite em
direção aos \omterm{def:g3:continents-and-oceans:poles}{polos}: os \omterm{def:g3:continents-and-oceans:poles}{polos} ficariam infinitamente longe do \omterm{def:g3:continents-and-oceans:equator}{equador}.
\end{solution}

\begin{solution}{exo:g9:maps-of-the-earth:6}
(a) \qty{50}{m}. (b) O cume está a \qty{544}{m}; a curva mestra em volta
dele é a linha de \qty{500}{m}. (c) No oeste: um \omterm{def:g1:landscapes-around-us:water}{rio} corre morro abaixo,
então o seu fundo é mais baixo \omterm{def:g1:landscapes-around-us:water}{rio} abaixo, a leste. (d) Mais ou menos
entre as linhas de 300 e de \qty{400}{m}: ali a encosta é mais íngreme,
mais suave tanto perto do fundo do \omterm{def:g1:landscapes-around-us:landform}{vale} quanto perto do cume.
\end{solution}

\begin{solution}{exo:g9:maps-of-the-earth:7}
(a) C a cerca de \qty{400}{m}, D a cerca de \qty{530}{m}. (b) Cerca de
\qty{135}{m}, a cerca de \qty{1.3}{km} de C, no fundo do \omterm{def:g1:landscapes-around-us:landform}{vale}. (c) Cerca
de 0.3 a \qty{0.4}{km}. (d) As alturas são pequenas em comparação com as
distâncias: sem esticar, os morros pareceriam quase planos.
\end{solution}

\begin{solution}{exo:g9:maps-of-the-earth:8}
(a) $2 \times 25\,000 = \num{50000}$~cm $= \qty{500}{m}$. (b) $300 - 200 =
\qty{100}{m}$. (c) $100 \div 500 = 0.2$, uma declividade de \qty{20}{\%}.
(d) Mais suave: \qty{4}{cm} são \qty{1000}{m}, então a declividade é
$100 \div 1000 = \qty{10}{\%}$, metade.
\end{solution}

\begin{solution}{exo:g9:maps-of-the-earth:9}
(a) $14 \times 18 = 252$ órbitas. (b) $14 \times 103 = 1442$ minutos,
cerca de \qty{24}{h}: um dia. (c) $5.5 \div 1 = 5.5$ vezes. (d) Por
exemplo: uma \omterm{def:g8:ice-on-earth:glacier}{geleira} que recua, uma floresta derrubada, uma cidade que
cresce, um \omterm{def:g1:landscapes-around-us:water}{litoral} ou um \omterm{def:g1:landscapes-around-us:water}{rio} que se desloca.
\end{solution}

\begin{solution}{exo:g9:maps-of-the-earth:10}
(a) $20\,200 \div 300\,000 \approx 0.067$~s $= 6.7 \times 10^{-2}$~s.
(b) $300\,000 \times 10^{-3} = \qty{300}{km}$. (c) Um erro de tempo
minúsculo causa um erro de distância enorme: os relógios têm de ser
extremamente precisos.
\end{solution}

\begin{solution}{exo:g9:maps-of-the-earth:11}
(a) Distância até A: $\sqrt{3.6^2 + 4.8^2} = \sqrt{12.96 + 23.04} =
\sqrt{36} = 6$; distância até B: $\sqrt{6.4^2 + 4.8^2} = \sqrt{40.96 +
23.04} = \sqrt{64} = 8$. (b) $(3.6, -4.8)$, a imagem espelhada em relação
à reta AB. (c) Distância de C até $(3.6, 4.8)$: $\sqrt{3.6^2 + 5.2^2} =
\sqrt{40} \approx 6.3$; até $(3.6, -4.8)$: $\sqrt{3.6^2 + 14.8^2} =
\sqrt{232} \approx 15.2$. O receptor está em $(3.6, 4.8)$. (d) Uma
distância dá um círculo (uma esfera no espaço), duas dão dois pontos
possíveis; é preciso outra distância para escolher.
\end{solution}

\begin{solution}{exo:g9:maps-of-the-earth:12}
(a) $29 \div 2.1 \approx 14$ vezes. (b) $2.1 \times 10.5 \approx
22$~milhões de~\unit{km^2}. (c) A Groenlândia parece cobrir cerca de
22~milhões de \unit{km^2}, perto dos 29~milhões da África: no \omterm{def:g2:drawing-to-map:map}{mapa} de
parede de Mercator, as duas parecem mais ou menos do mesmo tamanho,
embora a África seja 14 vezes maior. (d) Uma \omterm{def:g9:maps-of-the-earth:projection}{projeção equivalente}.
\end{solution}

\begin{solution}{pb:g9:maps-of-the-earth:1}
\textbf{1.} $50\,000$~cm $= \qty{500}{m}$.

\textbf{2.} $3.6 \times 500 = \qty{1800}{m} = \qty{1.8}{km}$.

\textbf{3.} B--W: $2.5 \times 500 = \qty{1250}{m}$; W--T: $2.35 \times 500
\approx \qty{1180}{m}$.

\textbf{4.} Para o nordeste.

\textbf{5.} \qty{50}{m}.

\textbf{6.} S--B desce: $296 - 143 = \qty{153}{m}$.

\textbf{7.} B--W: $335 - 143 = \qty{192}{m}$; W--T: $544 - 335 =
\qty{209}{m}$.

\textbf{8.} $209 \div 1180 \approx 0.18$, cerca de \qty{18}{\%}.

\textbf{9.} Uma linha reta de S até T atravessaria o \omterm{def:g1:landscapes-around-us:water}{rio} onde não há
ponte; a trilha desce até a ponte em B.

\textbf{10.} $1.8 + 1.25 + 1.18 \approx \qty{4.2}{km}$ a \qty{4}{km} por
hora: cerca de 63 minutos; subida $192 + 209 = \qty{401}{m}$: cerca de 40
minutos a mais. Ao todo, cerca de 1~h~43~min.

\textbf{11.} De T ela desce até B, depois sobe de B até S: $296 -
143 = \qty{153}{m}$.

\textbf{12.} Comprimento: $1.8 + 1.25 + 1.18 \approx \textbf{4.2 km}$;
subida total: $192 + 209 = \textbf{401 m}$.
\end{solution}
